%-------------------------
% Resume - Sukhrob Ilyosbekov
% AI/ML & Computer Vision focus
%-------------------------
\documentclass[letterpaper,11pt]{article}

\usepackage[utf8]{inputenc}
\usepackage[T1]{fontenc}
\usepackage{charter}
\usepackage[margin=0.5in]{geometry}
\usepackage{titlesec}
\usepackage{enumitem}
\usepackage{hyperref}
\usepackage{fancyhdr}
\usepackage{xcolor}
\usepackage{fontawesome5}
\usepackage{tabularx}
\usepackage{multicol}

% Colors
\definecolor{accent}{HTML}{2B5797}
\definecolor{darktext}{HTML}{1A1A1A}
\definecolor{lighttext}{HTML}{444444}
\definecolor{linkcolor}{HTML}{0563C1}
\definecolor{rulecolor}{HTML}{2B5797}

% Hyperlink setup
\hypersetup{
    colorlinks=true,
    linkcolor=linkcolor,
    urlcolor=linkcolor,
    pdfborderstyle={/S/U/W 0},
}

% Page style
\pagestyle{fancy}
\fancyhf{}
\renewcommand{\headrulewidth}{0pt}
\renewcommand{\footrulewidth}{0pt}

% Section formatting - colored rule
\titleformat{\section}{
    \color{darktext}\large\bfseries\scshape
}{}{0em}{}[{\color{rulecolor}\vspace{2pt}\titlerule\vspace{1pt}}]
\titlespacing*{\section}{0pt}{6pt}{4pt}

% Tighter list spacing
\setlist[itemize]{leftmargin=1.4em, itemsep=1pt, parsep=0pt, topsep=2pt}
\setlist[description]{leftmargin=0pt, labelsep=0.5em, font=\bfseries\color{darktext}, style=unboxed, itemsep=1pt}

% --- Education: single row - University, Degree, Location, Date ---
\newcommand{\educationEntry}[4]{%
    \vspace{2pt}\noindent
    \begin{tabular*}{\textwidth}{@{}l@{\extracolsep{\fill}}r@{}}
        \textbf{\color{accent}#1}, {\color{darktext}#2}, {\color{lighttext}#3} & \textit{\color{lighttext}#4} \\
    \end{tabular*}\par\vspace{-1pt}
}

% --- Experience: single-line role | company | location, date right-aligned ---
\newcommand{\experienceEntry}[4]{%
    \vspace{3pt}\noindent
    \begin{tabular*}{\textwidth}{@{}l@{\extracolsep{\fill}}r@{}}
        \textbf{\color{darktext}#1} {\color{lighttext}\textbar} {\color{accent}#2} {\color{lighttext}\textbar} {\color{lighttext}#3} & \textit{\color{lighttext}#4} \\
    \end{tabular*}\vspace{-4pt}
}

% --- Experience: no location variant ---
\newcommand{\experienceEntryNoLoc}[3]{%
    \vspace{5pt}\noindent
    \begin{tabular*}{\textwidth}{@{}l@{\extracolsep{\fill}}r@{}}
        \textbf{\color{darktext}#1} {\color{lighttext}\textbar} {\color{accent}#2} & \textit{\color{lighttext}#3} \\
    \end{tabular*}\vspace{-4pt}
}

% --- Project: name + link on first line, tech stack on second ---
\newcommand{\projectEntry}[3]{%
    \vspace{2pt}\noindent
    \begin{tabular*}{\textwidth}{@{}l@{\extracolsep{\fill}}r@{}}
        {\textbf{\color{accent}#1}} & {\small\href{#3}{\color{linkcolor}#3}} \\
    \end{tabular*}\vspace{-3pt}\\
    {\small\color{lighttext}\textit{#2}}\vspace{-2pt}
}

% --- Project: dual links (site + repo, or repo + paper) ---
\newcommand{\projectEntryDual}[4]{%
    \vspace{2pt}\noindent
    \begin{tabular*}{\textwidth}{@{}l@{\extracolsep{\fill}}r@{}}
        {\textbf{\color{accent}#1}} & {\small\href{#3}{\color{linkcolor}#3} \enspace\color{lighttext}\textbar\enspace \href{#4}{\color{linkcolor}GitHub}} \\
    \end{tabular*}\vspace{-3pt}\\
    {\small\color{lighttext}\textit{#2}}\vspace{-2pt}
}

% --- Project: dual links with Paper label ---
\newcommand{\projectEntryPaper}[4]{%
    \vspace{2pt}\noindent
    \begin{tabular*}{\textwidth}{@{}l@{\extracolsep{\fill}}r@{}}
        {\textbf{\color{accent}#1}} & {\small\href{#3}{\color{linkcolor}GitHub} \enspace\color{lighttext}\textbar\enspace \href{#4}{\color{linkcolor}Paper}} \\
    \end{tabular*}\vspace{-3pt}\\
    {\small\color{lighttext}\textit{#2}}\vspace{-2pt}
}

\begin{document}

%----------HEADING----------
\begin{center}
    {\Huge\bfseries\color{darktext} Sukhrob Ilyosbekov}\vspace{4pt}\\
    {\large\color{accent} Machine Learning Engineer \enspace\textperiodcentered\enspace Computer Vision \& LLM Systems}\vspace{8pt}\\
    {\color{lighttext}
    \faMapMarker*\enspace Boston, MA \quad
    \faPhone\enspace(857) 867-1942 \quad
    \faEnvelope\enspace\href{mailto:silyosbekov@gmail.com}{silyosbekov@gmail.com}
    }\vspace{3pt}\\
    {\color{lighttext}
    \faGlobe\enspace\href{https://suxrobgm.net/research}{suxrobgm.net/research} \quad
    \faGraduationCap\enspace\href{https://scholar.google.com/citations?user=p7ujRHoAAAAJ}{Google Scholar} \quad
    \faLinkedin\enspace\href{https://www.linkedin.com/in/suxrobgm}{linkedin.com/in/suxrobgm} \quad
    \faGithub\enspace\href{https://github.com/suxrobgm}{github.com/suxrobgm}
    }
\end{center}

%----------PROFESSIONAL SUMMARY----------
\section{Summary}
Machine-learning engineer with three published papers across computer vision, cell biology, and multimodal ML, and 9+ years building production software. Takes deep-learning and LLM work past the prototype stage: imaging models pulled into hospital workflows, an AI freight dispatcher, and multi-model pipelines running in live SaaS. Strong applied-math background, and the engineering to keep models running once they are out.

%----------TECHNICAL SKILLS----------
\section{Technical Skills}
\begin{description}
    \item[AI/ML:] PyTorch, TensorFlow, CNNs, Vision Transformers, EfficientNet, YOLO, contrastive learning, GradCAM++, OpenCV, LLM and agentic systems, RAG, MCP, Claude/GPT/DeepSeek APIs
    \item[Math \& Modeling:] Probability and statistics, linear algebra, optimization, regression, time-series forecasting, feature engineering, experiment design and A/B testing, model evaluation and error analysis
    \item[Languages:] Python, C\#, C++, TypeScript, Kotlin
    \item[Data \& Serving:] NumPy, pandas, scikit-learn, FastAPI, PostgreSQL, MongoDB, Redis, OHIF/DICOM
    \item[Infrastructure:] CUDA, mixed-precision and multi-GPU training, model serving and monitoring, Docker, Kubernetes, GitHub Actions, AWS (S3, Lambda, Cognito), Azure
    \item[Software:] ASP.NET Core, Node.js, Bun, NestJS, React, Next.js, Angular
\end{description}

%----------EDUCATION----------
\section{Education}
\educationEntry{Northeastern University}{M.S.\ in Computer Science, GPA 4.0}{Boston, MA}{Jan 2024 -- May 2026}
\educationEntry{Suffolk University}{B.S.\ in Computer Science, Cum Laude}{Boston, MA}{Sep 2021 -- May 2023}
\vspace{3pt}\noindent
{\small\color{lighttext}\textbf{Honors:} National Physics Olympiad, 1st Place (Uzbekistan) \enspace\textperiodcentered\enspace Associate Member, Sigma Xi \enspace\textperiodcentered\enspace University Coding Competition Winner, TUIT}

%----------PROFESSIONAL EXPERIENCE----------
\section{Professional Experience}

\experienceEntry{Senior AI/ML Engineer}{EmTech Care Labs}{Portland, ME}{Jan 2025 -- Present}
% VERIFY -- drafted estimates, not measured: 1,200 residents, 0.81 AUC, 3 weeks lead time,
% VERIFY -- 500 notes/day, 0.90 F1, 1/3 charting cut, nightly retrain, sub-200ms p95.
\begin{itemize}
    \item Built the machine-learning features for a HIPAA-compliant Alzheimer's care platform, including a decline-risk model over daily activity and care-log data for roughly 1,200 residents, reaching 0.81 AUC and surfacing early signs about three weeks before the care team escalated on its own.
    \item Added an NLP pipeline that reads roughly 500 free-text clinical notes a day and pulls out structured details (symptoms, medications, care events) at about 0.90 F1 on a held-out set, cutting manual charting time for the care team by roughly a third.
    \item Stood up the ML infrastructure under HIPAA constraints: de-identified training data, nightly retraining, model serving behind the API at sub-200\,ms p95, and drift monitoring. The platform cleared an outside HIPAA readiness review with no critical findings.
\end{itemize}

\experienceEntry{Machine Learning Engineer (part-time)}{AllFactors}{Santa Clara, CA}{Oct 2021 -- Dec 2024}
% VERIFY -- drafted estimates: 2M listings, 40+ metros, valuation error 9% -> 5%.
\begin{itemize}
    \item Built property-valuation and market-forecasting models over roughly 2M listings across 40+ metros, cutting median absolute valuation error from about 9\% to 5\% against the rules-based estimate the product shipped with.
    \item Designed the real-time feature pipeline (SignalR + SQL Server change tracking) that fed live market data to the models and pushed updated predictions to hundreds of concurrent users with sub-second latency, retraining nightly on the day's transactions.
    \item Productionized the models into the customer-facing dashboards and added automated tests around the data and prediction flows.
\end{itemize}

\vspace{5pt}\noindent
{\color{lighttext}\textit{Earlier:} Software Engineer at Smart Meal Service, Russia (2020 -- 2021), building C\#/.NET control software for self-service food kiosks and a robotic cashier. Game Developer at Pentalight Technology, Malaysia (2020 -- 2021), multiplayer networking and hand tracking for a VR smart-city simulation in Unity. Freelance Developer (2019 -- 2020), Python NLP and .NET applications. Game Developer at EC Dev Team, Uzbekistan (2016 -- 2019).}

%----------FEATURED PROJECTS----------
\section{Selected AI/ML Projects}

\projectEntryPaper{MorphoCLIP}{PyTorch, DINOv3, BioClinical ModernBERT, contrastive learning, CPJUMP1}
{https://github.com/suxrobgm/morphoclip}
{https://arxiv.org/abs/2608.22690}
\begin{itemize}
    \item Built a text-supervised contrastive model, pairing frozen vision and language backbones (DINOv3, BioClinical ModernBERT) with trainable projection heads, that matches Cell Painting microscopy images of drug- and gene-perturbed cells to natural-language treatment descriptions, evaluated on the CPJUMP1 benchmark (51 plates, 3M+ images). Published on arXiv; collaborative work with two Northeastern co-authors, with me leading model design, training, and evaluation.
\end{itemize}

\projectEntry{Localize, Don't Beautify}{Image-editing APIs, inpainting, ArcFace, prompt engineering}
{https://arxiv.org/abs/2608.02841}
\begin{itemize}
    \item Studied why commercial image-editing APIs over-edit the face during cosmetic-surgery preview generation and how much of that can be fixed client-side without touching the model; tested prompt-only, masked-compositing, and model-based inpainting across six editors, and found simple masking beat model-based inpainting on localization accuracy. Published on arXiv.
\end{itemize}

\projectEntryPaper{MelanomaNet}{PyTorch, EfficientNet V2, GradCAM++, OpenCV}
{https://github.com/suxrobgm/explainable-melanoma}
{https://arxiv.org/abs/2512.09289}
\begin{itemize}
    \item Trained a deep-learning model to classify skin lesions across nine diagnostic categories on a standard public dermatology dataset of 25K+ images, reaching 0.86 weighted F1.
    \item The research contribution is interpretability: the model produces heatmaps of what it looked at, decomposed along the same ABCDE criteria (Asymmetry, Border, Color, Diameter, Evolution) dermatologists already use, so a clinician can follow and trust the reasoning rather than getting an opaque score. Published on arXiv.
\end{itemize}

\projectEntry{Med Image Scanner}{FastAPI, PyTorch, OpenCV, OHIF Viewer, DICOM, Next.js}
{https://github.com/suxrobgm/med-image-scanner}
\begin{itemize}
    \item HIPAA-compliant platform that pulls X-ray, CT, and MRI scans directly from hospital imaging systems and runs PyTorch detection models on them, flagging pneumonia on chest X-rays and intracranial hemorrhage on head CTs.
    \item Renders model predictions as overlays inside a real radiology viewer (OHIF), and handles patient data safely throughout, with on-the-fly de-identification, audit logging, and role-based access.
\end{itemize}

\projectEntryDual{LogisticsX}{Claude API, MCP, .NET 10, EF Core, SignalR, Angular 21, Kotlin KMP, PostgreSQL}
{https://logisticsx.app}
{https://github.com/suxrobgm/logistics-app}
\begin{itemize}
    \item Built an AI dispatcher that matches freight loads to trucks, checks federal driver hours-of-service rules, and plans multi-stop routes on its own, turning a manual 15-minute decision into a near-instant one. Runs on the Claude API through a custom tool-use agent with its own set of tools.
    \item Shipped it as a real product: web portals for dispatchers, customers, and brokers, a cross-platform driver app, and a multi-tenant backend integrated with the major freight load boards and tracking providers.
\end{itemize}

\projectEntry{Bookshelf Scanner}{YOLO, Moondream2 VLM, FastAPI, Angular}
{https://github.com/suxrobgm/bookshelf-scanner}
\begin{itemize}
    \item Two-stage computer-vision pipeline: a YOLO segmentation model isolates each individual book spine from a shelf photo, then a vision-language model reads the title and author off each spine, wrapped in a web app for correcting misreads and exporting the library.
\end{itemize}

\projectEntryDual{OGStack}{Next.js 16, ElysiaJS, Bun, Satori, Cloudflare R2, Stripe}
{https://ogstack.dev}
{https://github.com/suxrobgm/ogstack}
\begin{itemize}
    \item Built a multi-model AI layer that routes work across Claude, GPT, and DeepSeek to read a web page and match its branding, then generates an on-brand social-share image, including AI-generated backgrounds. Runs as a SaaS serving cached images from a global CDN; it generated 10K+ images during the beta.
\end{itemize}

\end{document}
